% Options for packages loaded elsewhere
\PassOptionsToPackage{unicode}{hyperref}
\PassOptionsToPackage{hyphens}{url}
\PassOptionsToPackage{dvipsnames,svgnames,x11names}{xcolor}
\documentclass[
  spanish,
  10pt,
  a4paper,
]{article}
\usepackage{xcolor}
\usepackage[top=22mm,bottom=22mm,left=22mm,right=22mm]{geometry}
\usepackage{amsmath,amssymb}
\setcounter{secnumdepth}{-\maxdimen} % remove section numbering
\usepackage{iftex}
\ifPDFTeX
  \usepackage[T1]{fontenc}
  \usepackage[utf8]{inputenc}
  \usepackage{textcomp} % provide euro and other symbols
\else % if luatex or xetex
  \usepackage{unicode-math} % this also loads fontspec
  \defaultfontfeatures{Scale=MatchLowercase}
  \defaultfontfeatures[\rmfamily]{Ligatures=TeX,Scale=1}
\fi
\usepackage{lmodern}
\ifPDFTeX\else
  % xetex/luatex font selection
\fi
% Use upquote if available, for straight quotes in verbatim environments
\IfFileExists{upquote.sty}{\usepackage{upquote}}{}
\IfFileExists{microtype.sty}{% use microtype if available
  \usepackage[]{microtype}
  \UseMicrotypeSet[protrusion]{basicmath} % disable protrusion for tt fonts
}{}
\usepackage{setspace}
\makeatletter
\@ifundefined{KOMAClassName}{% if non-KOMA class
  \IfFileExists{parskip.sty}{%
    \usepackage{parskip}
  }{% else
    \setlength{\parindent}{0pt}
    \setlength{\parskip}{6pt plus 2pt minus 1pt}}
}{% if KOMA class
  \KOMAoptions{parskip=half}}
\makeatother
\usepackage{longtable,booktabs,array}
\newcounter{none} % for unnumbered tables
\usepackage{calc} % for calculating minipage widths
% Correct order of tables after \paragraph or \subparagraph
\usepackage{etoolbox}
\makeatletter
\patchcmd\longtable{\par}{\if@noskipsec\mbox{}\fi\par}{}{}
\makeatother
% Allow footnotes in longtable head/foot
\IfFileExists{footnotehyper.sty}{\usepackage{footnotehyper}}{\usepackage{footnote}}
\makesavenoteenv{longtable}
\usepackage{graphicx}
\makeatletter
\newsavebox\pandoc@box
\newcommand*\pandocbounded[1]{% scales image to fit in text height/width
  \sbox\pandoc@box{#1}%
  \Gscale@div\@tempa{\textheight}{\dimexpr\ht\pandoc@box+\dp\pandoc@box\relax}%
  \Gscale@div\@tempb{\linewidth}{\wd\pandoc@box}%
  \ifdim\@tempb\p@<\@tempa\p@\let\@tempa\@tempb\fi% select the smaller of both
  \ifdim\@tempa\p@<\p@\scalebox{\@tempa}{\usebox\pandoc@box}%
  \else\usebox{\pandoc@box}%
  \fi%
}
% Set default figure placement to htbp
\def\fps@figure{htbp}
\makeatother
\ifLuaTeX
\usepackage[bidi=basic,shorthands=off]{babel}
\else
\usepackage[bidi=default,shorthands=off]{babel}
\fi
\ifLuaTeX
  \usepackage{selnolig} % disable illegal ligatures
\fi
\setlength{\emergencystretch}{3em} % prevent overfull lines
\providecommand{\tightlist}{%
  \setlength{\itemsep}{0pt}\setlength{\parskip}{0pt}}

\usepackage{amsmath,amssymb,mathtools}
\usepackage{microtype}
\usepackage{fancyhdr}
\usepackage{needspace}
\usepackage{etoolbox}
\usepackage{float}
\usepackage{caption}
\captionsetup{font=small,labelformat=empty,skip=7pt}
\floatplacement{figure}{H}
\definecolor{DocBlue}{HTML}{173F57}
\definecolor{DocGray}{HTML}{52616C}
\pagestyle{fancy}
\fancyhf{}
\fancyhead[L]{\small\sffamily\color{DocGray} pySNSPD | Desarrollo del modelo}
\fancyhead[R]{\small\sffamily\color{DocGray} Revisión 0.4}
\fancyfoot[L]{\small\sffamily\color{DocGray} 9 de septiembre de 2026}
\fancyfoot[R]{\small\sffamily\thepage}
\setlength{\headheight}{15pt}
\setlength{\emergencystretch}{2em}
\setlength{\parskip}{5pt plus 1pt}
\setlength{\parindent}{0pt}
\widowpenalty=10000
\clubpenalty=10000
\displaywidowpenalty=10000
\predisplaypenalty=10000
\renewcommand{\arraystretch}{1.16}
\AtBeginEnvironment{longtable}{\small}
\pretocmd{\section}{\Needspace{6\baselineskip}}{}{}
\pretocmd{\subsection}{\Needspace{5\baselineskip}}{}{}
\AtBeginDocument{\hypersetup{colorlinks=true,linkcolor=DocBlue,urlcolor=DocBlue}}
\allowdisplaybreaks[1]
\usepackage{bookmark}
\IfFileExists{xurl.sty}{\usepackage{xurl}}{} % add URL line breaks if available
\urlstyle{same}
\hypersetup{
  pdftitle={E. Cuaderno de aprendizaje del modelo},
  pdflang={es},
  colorlinks=true,
  linkcolor={Maroon},
  filecolor={Maroon},
  citecolor={Blue},
  urlcolor={Blue},
  pdfcreator={LaTeX via pandoc}}

\title{E. Cuaderno de aprendizaje del modelo}
\usepackage{etoolbox}
\makeatletter
\providecommand{\subtitle}[1]{% add subtitle to \maketitle
  \apptocmd{\@title}{\par {\large #1 \par}}{}{}
}
\makeatother
\subtitle{Un concepto por paso · Documento 5 de 5}
\author{}
\date{9 de septiembre de 2026 · Revisión 0.4}

\begin{document}
\maketitle

{
\setcounter{tocdepth}{2}
\tableofcontents
}
\setstretch{1.07}
\newpage

\setlength{\parskip}{4pt plus 1pt}

\section{E.0. Ruta y registro de
aprendizaje}\label{e.0.-ruta-y-registro-de-aprendizaje}

Esta versión reduce los requisitos de entrada. Cada clase introduce sus
símbolos antes de calcular, distingue el objeto físico de su
representación matemática y trabaja un ejemplo pequeño. No hace falta
dominar teoría cuántica de campos, cálculo de operadores ni DFPT para
seguir el primer paso del modelo. El objetivo es comprender qué hace una
fórmula de A--D y qué información necesita.

No se recibieron respuestas en el Markdown de 0.3. Por tanto, ninguna
clase está aprobada. Una clase sin respuesta y sin otro motivo explícito
se interpreta desde ahora como una señal de dificultad excesiva:
corresponde revisar su enseñanza, no asignar una nota de reprobación.
Cuando haya clases cerradas, las siguientes partirán de esos conceptos
demostrados; por ahora las conexiones de la tabla son una ruta
propuesta, no conocimientos que se dan por adquiridos.

{\def\LTcaptype{none} % do not increment counter
\begin{longtable}[]{@{}
  >{\raggedright\arraybackslash}p{(\linewidth - 6\tabcolsep) * \real{0.2500}}
  >{\raggedright\arraybackslash}p{(\linewidth - 6\tabcolsep) * \real{0.2500}}
  >{\raggedright\arraybackslash}p{(\linewidth - 6\tabcolsep) * \real{0.2500}}
  >{\raggedright\arraybackslash}p{(\linewidth - 6\tabcolsep) * \real{0.2500}}@{}}
\toprule\noalign{}
\begin{minipage}[b]{\linewidth}\raggedright
ID estable
\end{minipage} & \begin{minipage}[b]{\linewidth}\raggedright
Clase de esta versión
\end{minipage} & \begin{minipage}[b]{\linewidth}\raggedright
Estado
\end{minipage} & \begin{minipage}[b]{\linewidth}\raggedright
Punto de apoyo
\end{minipage} \\
\midrule\noalign{}
\endhead
\bottomrule\noalign{}
\endlastfoot
E01 & Partículas, excitaciones y escalas & ABIERTO & Variables
familiares de mecánica y circuitos \\
E04 & Electrones y huecos: elegir la referencia & ABIERTO & Carga y
portadores de un diodo \\
E06 & Ocupaciones y operadores: primera segunda cuantización & ABIERTO &
Columnas y multiplicación de matrices \\
E07 & Por qué se usa una columna de Nambu & ABIERTO & E04 y E06, sin
exigir matrices de cuatro componentes \\
E03 & Qué significa variar una energía funcional & ABIERTO & Derivada
ordinaria, regla del producto e integral \\
E05 & De dos masas a una DOS & ABIERTO & Modos normales de mecánica \\
E08 & Qué aporta DFPT al modelo del detector & ABIERTO & E05:
frecuencias, modos y conteo \\
E02 & Nambu y espín: dos índices distintos, clase de 0.3 & DESCARTADA &
Sustituida por E06 y E07; no se conserva su contenido \\
\end{longtable}
}

Los números E01, E03, etc. identifican clases a través de versiones;
E.1, E.2, etc. indican su orden de lectura actual. \textbf{DESCARTADA}
identifica exclusivamente una clase retirada y no es una evaluación del
aprendizaje. Las clases activas tienen los estados ABIERTO o CERRADO /
APROBADO.

Cada clase vale 10 puntos: 4 por un cálculo guiado, 3 por interpretar y
3 por detectar una confusión. Se cierra con al menos 8 puntos y la
condición esencial satisfecha. Los errores aritméticos menores reciben
crédito parcial. Las respuestas se escriben bajo su identificador. La
próxima versión conservará las respuestas y añadirá la pauta de las
actividades entregadas. Aquí se resuelven ejemplos distintos.

\subsection{E.0.1. La palabra «estado» necesita
apellido}\label{e.0.1.-la-palabra-estado-necesita-apellido}

{\def\LTcaptype{none} % do not increment counter
\begin{longtable}[]{@{}
  >{\raggedright\arraybackslash}p{(\linewidth - 4\tabcolsep) * \real{0.3333}}
  >{\raggedright\arraybackslash}p{(\linewidth - 4\tabcolsep) * \real{0.3333}}
  >{\raggedright\arraybackslash}p{(\linewidth - 4\tabcolsep) * \real{0.3333}}@{}}
\toprule\noalign{}
\begin{minipage}[b]{\linewidth}\raggedright
Expresión
\end{minipage} & \begin{minipage}[b]{\linewidth}\raggedright
Qué significa aquí
\end{minipage} & \begin{minipage}[b]{\linewidth}\raggedright
Ejemplo
\end{minipage} \\
\midrule\noalign{}
\endhead
\bottomrule\noalign{}
\endlastfoot
Estado mecánico clásico & Coordenadas y momentos que permiten continuar
la evolución & Posición y momento de dos masas \\
Vector de estado del circuito & Variables independientes de
almacenamiento elegidas para el circuito & Corriente de una inductancia
y tensión de un capacitor \\
Estado cuántico de ocupación & Vector que especifica las ocupaciones, o
una superposición de esos vectores & Un modo electrónico vacío u
ocupado \\
Estado electrónico de una partícula & Modo disponible antes de decir si
está ocupado & Un orbital con una orientación de espín definida \\
Estado químico & Configuración, enlaces o grado de oxidación, según la
pregunta & Átomo neutro e ion; no equivale a toda su función de onda \\
Estado del condensado & Campo de orden efectivo usado por el modelo &
Perfil complejo \(\Delta(\mathbf r,t)\) \\
\end{longtable}
}

El mismo sistema admite varias descripciones. Una corriente de circuito
no especifica cada ocupación electrónica; un listado de especies
elementales tampoco determina un pulso eléctrico.

\Needspace{12\baselineskip}

\section{E.1. Partículas, excitaciones y escalas ---
E01}\label{e.1.-partuxedculas-excitaciones-y-escalas-e01}

\textbf{Estado: ABIERTO. Meta:} distinguir qué se está describiendo
antes de interpretar masa, carga o espín.

\subsection{E.1.1. Introducción}\label{e.1.1.-introducciuxf3n}

En mecánica, una coordenada representa el desplazamiento de una masa. En
una descripción de campos, una variable se asigna a cada posición. Un
modo es una forma de variación de ese campo, del mismo modo que un modo
normal es una forma de vibración de una cadena. Una excitación cuántica
es un cambio permitido de energía y ocupación de esos modos. Esta
analogía introduce el lenguaje; no construye literalmente un electrón a
partir de resortes.

\begin{figure}
\centering
\includegraphics[width=1\linewidth,height=\textheight,keepaspectratio,alt={Figura E.1. Tres escalas de descripción. El vocabulario y las variables cambian con la pregunta. El fonón pertenece a la descripción de un sólido; la corriente pertenece a la descripción del dispositivo.}]{../figuras/E01_escalas.png}
\caption{Figura E.1. Tres escalas de descripción. El vocabulario y las
variables cambian con la pregunta. El fonón pertenece a la descripción
de un sólido; la corriente pertenece a la descripción del dispositivo.}
\end{figure}

Una \textbf{especie de partícula} identifica propiedades compartidas:
todos los electrones tienen la misma masa y carga. Un \textbf{estado
cuántico} indica cómo se encuentra o se ocupa el sistema descrito. Son
preguntas distintas. El espín es momento angular intrínseco y describe
cómo transforma el estado bajo rotaciones; no es una bolita girando.

El Modelo Estándar organiza campos elementales y sus interacciones
electromagnética, débil y fuerte. Los quarks y leptones son fermiones:
un estado electrónico o fermiónico completamente especificado admite
ocupación 0 o 1. Los bosones admiten otras ocupaciones. La tabla muestra
ejemplos, no el inventario completo.

{\def\LTcaptype{none} % do not increment counter
\begin{longtable}[]{@{}llll@{}}
\toprule\noalign{}
Familia & Ejemplos & Carga eléctrica, en unidades de \(e>0\) & Espín \\
\midrule\noalign{}
\endhead
\bottomrule\noalign{}
\endlastfoot
Quarks & up, down & \(+2/3\), \(-1/3\) & \(1/2\) \\
Leptones & electrón, neutrino & \(-1\), \(0\) & \(1/2\) \\
Bosones asociados a interacciones & fotón, \(W^+\) & \(0\), \(+1\) &
\(1\) \\
Bosón de Higgs & Higgs & \(0\) & \(0\) \\
\end{longtable}
}

Masa, carga y espín no especifican por sí solos el estado de movimiento,
la ocupación o las correlaciones. El Modelo Estándar tampoco incluye una
teoría cuántica completa de la gravedad. Un protón es compuesto; un
fonón es una excitación colectiva de un sólido, no otra partícula
elemental de esa tabla.

\subsection{E.1.2. Ejemplo resuelto}\label{e.1.2.-ejemplo-resuelto}

\textbf{Problema:} un protón contiene quarks de valencia up, up, down.
¿Qué se puede obtener sumando los datos de carga?

La suma es \(2/3+2/3-1/3=1\): carga \(+e\). Si se añade un electrón de
carga \(-e\), el conjunto tiene carga neta cero. Eso no determina la
masa, la energía de enlace ni el estado cuántico. La carga es aditiva;
la energía de un sistema ligado incluye interacciones. En particular,
sumar tres números de espín \(1/2\) tampoco determina el espín del
protón: hay que conocer cómo se combinan los momentos angulares.

\subsection{E.1.3. Actividad}\label{e.1.3.-actividad}

\Needspace{4\baselineskip}

\textbf{E01.1, 4 puntos.} Un neutrón tiene quarks de valencia up, down,
down. Calcular su carga con la tabla e indicar una propiedad que no
queda determinada por esa suma.

\textbf{Respuesta E01.1:} \emph{Escribir aquí.}

\Needspace{4\baselineskip}

\textbf{E01.2, 3 puntos.} Clasificar electrón, protón, fonón y corriente
de una inductancia como especie elemental, objeto compuesto, excitación
colectiva o variable de circuito.

\textbf{Respuesta E01.2:} \emph{Escribir aquí.}

\Needspace{4\baselineskip}

\textbf{E01.3, 3 puntos.} Corregir: «si dos electrones tienen igual
masa, carga y espín, están necesariamente en el mismo estado cuántico».

\textbf{Respuesta E01.3:} \emph{Escribir aquí.}

\textbf{Condición esencial:} separar especie, estado cuántico y variable
efectiva.

\textbf{Conexión con A--D:} A describe electrones próximos al nivel de
Fermi; B sigue ocupaciones de excitaciones; C utiliza el campo efectivo
del condensado; D reúne además variables de circuito. Fuentes de
contexto: \href{https://home.cern/science/physics/standard-model/}{CERN,
Modelo Estándar},
\href{https://pdg.lbl.gov/chris/museum_version/standard_model.html}{Particle
Data Group, clasificación} y
\href{https://www.damtp.cam.ac.uk/user/tong/qft/qfthtml/S0.html}{Tong,
introducción a campos}. La analogía y el ejercicio se elaboran aquí.

\Needspace{12\baselineskip}

\section{E.2. Electrones y huecos: elegir la referencia ---
E04}\label{e.2.-electrones-y-huecos-elegir-la-referencia-e04}

\textbf{Estado: ABIERTO. Meta:} contabilizar carga y energía sin
convertir una vacante en una nueva especie elemental.

\subsection{E.2.1. Introducción desde el
diodo}\label{e.2.1.-introducciuxf3n-desde-el-diodo}

En una banda casi llena de un semiconductor, seguir los pocos lugares
vacíos puede ser más sencillo que seguir todos los electrones. Esa
vacante se llama hueco. Si un electrón llena la vacante vecina, la
vacante cambia de posición en sentido contrario. La carga relativa de
retirar un electrón es positiva, porque se quitó una carga negativa.

En un diodo se distinguen electrones de conducción y huecos de valencia.
Su difusión y la acción del campo eléctrico determinan corrientes; en
equilibrio pueden compensarse. La analogía que se usará para el
superconductor es la \textbf{contabilidad respecto de una referencia
ocupada}, no el perfil de bandas completo del diodo.

\begin{figure}
\centering
\includegraphics[width=0.98\linewidth,height=\textheight,keepaspectratio,alt={Figura E.2. Izquierda: un salto electrónico mueve la vacante al revés. Derecha: se pueden crear excitaciones de energía positiva añadiendo por encima de la referencia o retirando por debajo.}]{../figuras/E04_referencia_hueco.png}
\caption{Figura E.2. Izquierda: un salto electrónico mueve la vacante al
revés. Derecha: se pueden crear excitaciones de energía positiva
añadiendo por encima de la referencia o retirando por debajo.}
\end{figure}

Antes de calcular, se definen \(\varepsilon\) como energía de un nivel
electrónico, \(\mu\) como potencial químico de referencia y \(\delta N\)
como cambio del número de electrones: \(+1\) al añadir, \(-1\) al
retirar. La carga relativa es \(\delta Q=-e\delta N\). La energía de
excitación respecto de esa referencia es el cambio de \(U-\mu N\), no
necesariamente el cambio de energía interna \(U\) aislada:

\[
\delta E_{\rm exc}=(\varepsilon-\mu)\delta N,
\qquad \delta Q=-e\delta N.
\tag{E.1}
\]

La forma es «energía por electrón medida desde la referencia» por
«cantidad añadida o retirada». Retirar un electrón con
\(\varepsilon<\mu\) multiplica dos signos negativos y cuesta energía
positiva. Un hueco en el sólido no es un positrón.

\subsection{E.2.2. Ejemplo resuelto}\label{e.2.2.-ejemplo-resuelto}

Añadir un electrón a \(\varepsilon-\mu=3\) meV da
\(\delta E_{\rm exc}=3\) meV y \(\delta Q=-e\). Retirar otro a
\(\varepsilon-\mu=-2\) meV da \((-2)(-1)=2\) meV y \(\delta Q=+e\). El
conjunto contiene 5 meV de excitaciones y carga neta cero. Neutralidad
no significa ausencia de energía.

Una cuasipartícula superconductora es una excitación del sistema
electrónico emparejado. Su descripción combina operaciones de añadir y
retirar electrones. E07 introducirá esa combinación después de definir
dichas operaciones en E06; aquí no se exige su álgebra.

\subsection{E.2.3. Actividad}\label{e.2.3.-actividad}

\Needspace{4\baselineskip}

\textbf{E04.1, 4 puntos.} Añadir un electrón a \(\mu+4\) meV y retirar
otro a \(\mu-1\) meV. Completar por separado \(\delta N\), \(\delta Q\)
y \(\delta E_{\rm exc}\) de cada operación y sus totales.

\textbf{Respuesta E04.1:} \emph{Escribir aquí.}

\Needspace{4\baselineskip}

\textbf{E04.2, 3 puntos.} En la figura, si la vacante va a la derecha,
¿hacia dónde saltó el electrón y hacia dónde apunta la corriente
convencional asociada a ese salto? Explicar usando su carga.

\textbf{Respuesta E04.2:} \emph{Escribir aquí.}

\Needspace{4\baselineskip}

\textbf{E04.3, 3 puntos.} Corregir: «una población de excitaciones
eléctricamente neutra debe tener energía cero».

\textbf{Respuesta E04.3:} \emph{Escribir aquí.}

\textbf{Condición esencial:} declarar la referencia y conservar los
signos de carga y energía.

\textbf{Conexión con A--D:} el conteo normal de A incluye electrones y
vacantes; B integra energías de excitación positivas; C y D conservan
energía y carga mediante balances distintos. Fuentes:
\href{https://web.mit.edu/6.012/www/}{MIT 6.012, portadores y unión
p--n} y \href{https://topocondmat.org/test/w1_topointro/0d.html}{TU
Delft, representación superconductora}. Los números del ejemplo no son
parámetros del detector.

\Needspace{12\baselineskip}

\section{E.3. Ocupaciones y operadores ---
E06}\label{e.3.-ocupaciones-y-operadores-e06}

\textbf{Estado: ABIERTO. Meta:} leer \(c\), \(c^\dagger\) y
\(c^\dagger c\) como operaciones concretas. No se requiere teoría de
campos relativista.

\subsection{E.3.1. Introducción: cambiar de coordenadas de
descripción}\label{e.3.1.-introducciuxf3n-cambiar-de-coordenadas-de-descripciuxf3n}

En lugar de preguntar «dónde está cada electrón numerado», la
representación de ocupaciones pregunta «qué modos electrónicos están
ocupados». A este lenguaje de muchos cuerpos se lo llama \textbf{segunda
cuantización}; no significa aplicar dos veces un procedimiento de
cuantización a la misma variable.

Se empieza con un solo modo fermiónico, incluyendo su espín. Los
símbolos \(|0\rangle\) y \(|1\rangle\), llamados \emph{kets}, serán aquí
columnas ordinarias: vacío y ocupado. Una superposición es una
combinación lineal de esas columnas. Los operadores son matrices que
actúan sobre ellas.

\[
|0\rangle=\begin{pmatrix}1\\0\end{pmatrix},\quad
|1\rangle=\begin{pmatrix}0\\1\end{pmatrix},\quad
c=\begin{pmatrix}0&1\\0&0\end{pmatrix},\quad
c^\dagger=\begin{pmatrix}0&0\\1&0\end{pmatrix}.
\tag{E.2}
\]

La daga \(\dagger\) significa transponer y conjugar complejamente una
matriz. Como aquí las entradas son reales, basta transponer. El operador
\(c\) retira una ocupación y \(c^\dagger\) la añade. La forma de las
matrices se escoge precisamente para esas acciones:

\[
c|1\rangle=|0\rangle,\quad c|0\rangle=0,
\qquad c^\dagger|0\rangle=|1\rangle,\quad c^\dagger|1\rangle=0.
\tag{E.3}
\]

El \(0\) de la derecha es el \textbf{vector cero}, no el ket de vacío
\(|0\rangle\). Indica que la operación no produce un estado permitido.
En particular, no se añade un segundo fermión al mismo modo
completamente especificado.

\begin{figure}
\centering
\includegraphics[width=0.97\linewidth,height=\textheight,keepaspectratio,alt={Figura E.3. Acciones de añadir y retirar en un modo fermiónico. «Vacío» sigue siendo un estado de ocupación representado por un vector no nulo.}]{../figuras/E06_ocupacion_operadores.png}
\caption{Figura E.3. Acciones de añadir y retirar en un modo fermiónico.
«Vacío» sigue siendo un estado de ocupación representado por un vector
no nulo.}
\end{figure}

\subsection{E.3.2. Ejemplo resuelto: contar sin retirar
definitivamente}\label{e.3.2.-ejemplo-resuelto-contar-sin-retirar-definitivamente}

El operador número \(\hat n\) se define como \(c^\dagger c\). La matriz
de la derecha actúa primero:

\[
\hat n=c^\dagger c=\begin{pmatrix}0&0\\0&1\end{pmatrix},
\qquad \hat n|0\rangle=0,\quad \hat n|1\rangle=|1\rangle.
\tag{E.4}
\]

En el ket ocupado, primero se retira y luego se repone: el resultado es
el mismo ket multiplicado por 1. En el vacío, retirar da el vector cero.
Por eso los valores de ocupación son 0 y 1.

Al multiplicar en el otro orden se obtiene
\(cc^\dagger=\operatorname{diag}(1,0)\). Por tanto
\(cc^\dagger+c^\dagger c=I\), donde \(I\) es la matriz identidad. La
suma \(AB+BA\) se llama \textbf{anticomutador}. Este nombre permite
buscar la regla después; aquí se la ha comprobado con matrices. No se
necesita todavía el álgebra de muchos modos.

\subsection{E.3.3. Actividad}\label{e.3.3.-actividad}

\Needspace{4\baselineskip}

\textbf{E06.1, 4 puntos.} Multiplicar explícitamente
\(c^\dagger c|1\rangle\) y \(cc^\dagger|1\rangle\). Indicar en qué paso
importa el orden.

\textbf{Respuesta E06.1:} \emph{Escribir aquí.}

\Needspace{4\baselineskip}

\textbf{E06.2, 3 puntos.} Explicar la diferencia entre el ket
\(|0\rangle\) y el vector cero de E.3.

\textbf{Respuesta E06.2:} \emph{Escribir aquí.}

\Needspace{4\baselineskip}

\textbf{E06.3, 3 puntos.} Hay dos modos distintos, uno de espín arriba y
otro abajo. ¿La prohibición de dos fermiones en el mismo modo impide
ocupar uno en cada modo? Identificar qué etiqueta cambió.

\textbf{Respuesta E06.3:} \emph{Escribir aquí.}

\textbf{Condición esencial:} distinguir estado de ocupación, operador y
resultado de una operación.

\textbf{Conexión con A--D:} los operadores preparan la notación de Nambu
de A; las probabilidades de ocupación de B son promedios, no matrices
\(c\). C y D usan campos y variables de circuito derivados, no kets
individuales. Fuente de contexto:
\href{https://www.damtp.cam.ac.uk/user/tong/qft.htm}{Tong, cuantización
y representación de ocupaciones}. Las matrices y sus productos se
desarrollan aquí.

\Needspace{12\baselineskip}

\section{E.4. Por qué se usa una columna de Nambu ---
E07}\label{e.4.-por-quuxe9-se-usa-una-columna-de-nambu-e07}

\textbf{Estado: ABIERTO. Meta:} entender el motivo de juntar una
operación de retirar con otra de añadir. Requiere las operaciones de
E06.

\subsection{E.4.1. Introducción: el emparejamiento mezcla
operaciones}\label{e.4.1.-introducciuxf3n-el-emparejamiento-mezcla-operaciones}

Se eligen dos modos electrónicos relacionados: por ejemplo, momentos
opuestos y espines opuestos. Se los etiqueta 1 y 2; \(c_1\) y \(c_2\)
retiran sus ocupaciones. En una descripción efectiva del emparejamiento
aparecen términos que añaden o retiran una pareja,
\(\Delta c_1^\dagger c_2^\dagger\) y \(\Delta^*c_2c_1\). El orden
escrito es la convención de este ejemplo. El número complejo \(\Delta\)
mide el emparejamiento; su conjugado \(\Delta^*\) acompaña al proceso
inverso para que la energía sea real.

La energía normal de cada modo, medida desde el potencial químico, se
llama \(\xi\). Para representar juntos los términos normales y de
pareja, se define la \textbf{columna de operadores de Nambu}

\[
\Psi=\begin{pmatrix}c_1\\c_2^\dagger\end{pmatrix},\qquad
H_{\rm BdG}=\begin{pmatrix}\xi&\Delta\\\Delta^*&-\xi\end{pmatrix}.
\tag{E.5}
\]

BdG abrevia Bogoliubov--de Gennes. \(H_{\rm BdG}\) es una matriz de
coeficientes de energía. \textbf{\(\Psi\) contiene operadores; no es el
ket \(|0\rangle\) ni una lista de dos electrones nuevos.} La forma de la
matriz puede comprobarse expandiendo el producto siguiente.

\begin{figure}
\centering
\includegraphics[width=1\linewidth,height=\textheight,keepaspectratio,alt={Figura E.4. Multiplicar una matriz por una columna de operadores combina operaciones de retirar y añadir. El emparejamiento aparece fuera de la diagonal.}]{../figuras/E07_columna_nambu.png}
\caption{Figura E.4. Multiplicar una matriz por una columna de
operadores combina operaciones de retirar y añadir. El emparejamiento
aparece fuera de la diagonal.}
\end{figure}

\subsection{E.4.2. Ejemplo resuelto: de la matriz al contenido
físico}\label{e.4.2.-ejemplo-resuelto-de-la-matriz-al-contenido-fuxedsico}

Para esta pareja de modos, sin sumar otra vez el bloque equivalente,
multiplicar fila, matriz y columna da

\[
\begin{aligned}
\Psi^\dagger H_{\rm BdG}\Psi
={}&\xi c_1^\dagger c_1-\xi c_2c_2^\dagger
+\Delta c_1^\dagger c_2^\dagger+\Delta^*c_2c_1\\
={}&\xi(\hat n_1+\hat n_2-1)
+\Delta c_1^\dagger c_2^\dagger+\Delta^*c_2c_1.
\end{aligned}\tag{E.6}
\]

La segunda línea usa sólo \(c_2c_2^\dagger=1-\hat n_2\), que se practicó
en E06. El signo negativo de la diagonal no afirma que un segundo
electrón tenga energía normal \(-\xi\). Aparece porque se escribió
creación en esa entrada. Sumando la constante \(\xi\) se recupera la
energía normal \(\xi(\hat n_1+\hat n_2)\) más los términos de pareja.

Con \(\xi=2\) y \(\Delta=1\) en una misma unidad de energía, los valores
propios se obtienen de

\[
\det(H_{\rm BdG}-EI)=(2-E)(-2-E)-1=E^2-5=0,
\qquad E=\pm\sqrt5.
\tag{E.7}
\]

Un valor propio caracteriza un canal de la representación BdG. Las ramas
positivas y negativas están relacionadas por su construcción
electrón--hueco; no se cuentan como dos especies independientes. A
utiliza además índices de espín y trazas. Ese paso sólo se abordará
después de comprobar que E.5--E.6 se entienden; no hace falta
reconstruir ahora sus cuatro componentes.

\subsection{E.4.3. Actividad}\label{e.4.3.-actividad}

\Needspace{4\baselineskip}

\textbf{E07.1, 4 puntos.} Poner \(\Delta=0\) en E.6 y usar la constante
\(+\xi\). Mostrar que quedan las energías normales de los dos modos con
el mismo signo \(+\xi\).

\textbf{Respuesta E07.1:} \emph{Escribir aquí.}

\Needspace{4\baselineskip}

\textbf{E07.2, 3 puntos.} Señalar en E.5 qué entrada describe retirar
una ocupación y cuál añadirla. Explicar por qué ninguna entrada es una
probabilidad \(f(E)\).

\textbf{Respuesta E07.2:} \emph{Escribir aquí.}

\Needspace{4\baselineskip}

\textbf{E07.3, 3 puntos.} Corregir: «los dos signos de E.7 demuestran
que aparecieron dos tipos nuevos de partículas elementales».

\textbf{Respuesta E07.3:} \emph{Escribir aquí.}

\textbf{Condición esencial:} explicar el origen de la columna y el signo
diagonal mediante el producto, sin duplicar electrones físicos.

\textbf{Conexión con A--D:} A organiza el propagador electrónico con
esta clase de índices; B usa el espectro y la coherencia resultantes; C
evoluciona el emparejamiento efectivo; D conserva la misma convención al
formar fuerzas y corrientes. Fuente:
\href{https://topocondmat.org/test/w1_topointro/0d.html}{TU Delft,
Hamiltoniano BdG y redundancia}. Aquí se escogió un bloque de pareja
concreto; el factor \(1/2\) de una representación completa duplicada no
se traslada automáticamente a este bloque.

\Needspace{12\baselineskip}

\section{E.5. Qué significa variar una energía funcional ---
E03}\label{e.5.-quuxe9-significa-variar-una-energuxeda-funcional-e03}

\textbf{Estado: ABIERTO. Meta:} reconstruir una primera variación y
entender la envolvente de A. No se exige resolver un problema completo
de estabilidad ni calcular segundas variaciones.

\subsection{E.5.1. Introducción: construir primero una curva
ordinaria}\label{e.5.1.-introducciuxf3n-construir-primero-una-curva-ordinaria}

Una función \(F(y)\) recibe un número. Un funcional \(\mathcal F[y]\)
recibe un perfil completo \(y(x)\) y devuelve un número, por ejemplo una
energía. Los corchetes sólo recuerdan qué tipo de entrada se utiliza. La
coordenada \(x\) recorre el espacio; no representa tiempo en esta clase.

Se elige otro perfil \(\eta(x)\), llamado \textbf{dirección de
perturbación}, y un número pequeño \(\epsilon\) que fija su tamaño. El
perfil ensayado es \(y_\epsilon(x)=y(x)+\epsilon\eta(x)\). Una vez
fijados \(y\) y \(\eta\), evaluar su energía produce la función
ordinaria \(J(\epsilon)=\mathcal F[y_\epsilon]\) de una sola variable.

\begin{figure}
\centering
\includegraphics[width=0.96\linewidth,height=\textheight,keepaspectratio,alt={Figura E.5. La variación se convierte en la pendiente de una curva ordinaria. En el ejemplo dibujado, el término cuadrático explica por qué la recta tangente y la curva se separan al alejarse de cero.}]{../figuras/E03_curva_variacion.png}
\caption{Figura E.5. La variación se convierte en la pendiente de una
curva ordinaria. En el ejemplo dibujado, el término cuadrático explica
por qué la recta tangente y la curva se separan al alejarse de cero.}
\end{figure}

La derivada direccional de Gâteaux, también llamada primera variación en
esta aplicación, es

\[
D\mathcal F[y](\eta)
\equiv J'(0)
=\lim_{\epsilon\to0}
\frac{\mathcal F[y+\epsilon\eta]-\mathcal F[y]}{\epsilon}.
\tag{E.8}
\]

\(D\mathcal F[y](\eta)\) es un número que depende del perfil base y de
la dirección elegida. Cuando se puede escribir como
\(\int g(x)\eta(x)dx\), la función \(g(x)\) se denota
\(\delta\mathcal F/\delta y(x)\) y se llama densidad de derivada
funcional. \textbf{Ese símbolo no es un cociente de dos números
pequeños.} Los términos de borde deben tratarse antes de identificar
\(g\). Ésta es la interpretación necesaria de la antigua E.7 {[}MIT,
apéndice A, pp.~123--124{]}.

\subsection{E.5.2. Ejemplo resuelto: por qué quedan esos
términos}\label{e.5.2.-ejemplo-resuelto-por-quuxe9-quedan-esos-tuxe9rminos}

Se usa una energía adimensional sobre \(0\leq x\leq1\):

\[
\mathcal F[y]=\frac12\int_0^1[(y')^2+y^2]dx,
\qquad y'=\frac{dy}{dx}.
\tag{E.9}
\]

El primer cuadrado penaliza diferencias espaciales; el segundo, la
amplitud del perfil. Para hacer la variación se sustituye \textbf{en
ambos} cuadrados. Como \(\epsilon\) no depende de \(x\),
\((y+\epsilon\eta)'=y'+\epsilon\eta'\). La identidad
\((a+b)^2=a^2+2ab+b^2\) da exactamente

\[
\mathcal F[y+\epsilon\eta]-\mathcal F[y]
=\epsilon\int_0^1(y'\eta'+y\eta)dx
+\frac{\epsilon^2}{2}\int_0^1[(\eta')^2+\eta^2]dx.
\tag{E.10}
\]

No se ha descartado nada. Al dividir por \(\epsilon\), el primer término
ya no depende de él y el segundo conserva un factor \(\epsilon\). El
límite de E.8 elimina ese segundo término. Se retienen los términos
lineales \textbf{porque se está calculando una derivada en cero}, no
porque los demás sean físicamente inexistentes.

Para reemplazar \(\eta'\) por \(\eta\) se utiliza integración por
partes, que viene de integrar la regla del producto
\((y'\eta)'=y''\eta+y'\eta'\):

\[
D\mathcal F[y](\eta)
=[y'\eta]_0^1+\int_0^1(-y''+y)\eta\,dx.
\tag{E.11}
\]

Si los valores de \(y\) en los extremos están fijados, las
perturbaciones permitidas cumplen \(\eta(0)=\eta(1)=0\). Con esa
condición desaparece el término de borde y se identifica
\(\delta\mathcal F/\delta y=-y''+y\). Si los extremos pueden moverse,
ese término informa del trabajo de borde y se conserva hasta especificar
qué otra condición corresponde.

\textbf{Comprobación numérica del dibujo.} Para \(y=1\) y
\(\eta=\sin(\pi x)\), se obtiene \(D\mathcal F[y](\eta)=2/\pi\). La
curva exacta es \(J(\epsilon)=1/2+2\epsilon/\pi+(\pi^2+1)\epsilon^2/4\).
Su pendiente en cero vuelve a ser \(2/\pi\). El script verifica esa
pendiente con diferencias centrales.

\subsection{E.5.3. La envolvente, con una sola regla de la
cadena}\label{e.5.3.-la-envolvente-con-una-sola-regla-de-la-cadena}

En A hay variables internas que ya satisfacen su condición estacionaria.
Para ver qué cambia al eliminarlas, usar \(F(y,z)=y^2/2+(z-2y)^2/2\),
donde \(y,z\) son ahora números, no perfiles. La variable interna
estacionaria satisface \(F_z=z-2y=0\), así que \(z^*(y)=2y\). Definir
\(\bar F(y)=F(y,z^*(y))\) y aplicar la regla de la cadena:

\[
\frac{d\bar F}{dy}=F_y+F_z\frac{dz^*}{dy}=F_y,
\qquad F_z(y,z^*(y))=0.
\tag{E.12}
\]

Ésta es la forma estacionaria del \textbf{teorema de la envolvente},
para una rama diferenciable. La derivada interna no se elimina porque
\(z^*\) sea constante, sino porque multiplica una pendiente cero. Aquí
\(\bar F=y^2/2\) y la derivada vale \(y\). En A, resolver la variable
espectral interna no obliga a que la amplitud del condensado esté
también en equilibrio. Este paso y E.11 son el objetivo de la clase; el
análisis de curvaturas puede esperar.

\subsection{E.5.4. Actividad}\label{e.5.4.-actividad}

\Needspace{4\baselineskip}

\textbf{E03.1, 4 puntos.} Para
\(\mathcal G[y]=\frac12\int_0^1 y(x)^2dx\), expandir
\(\mathcal G[y+\epsilon\eta]-\mathcal G[y]\), dividir por \(\epsilon\) y
tomar el límite. No hay derivadas espaciales en este ejercicio.

\textbf{Respuesta E03.1:} \emph{Escribir aquí.}

\Needspace{4\baselineskip}

\textbf{E03.2, 3 puntos.} En E.11, explicar qué condición permite quitar
\([y'\eta]_0^1\). ¿Basta con decir «se integra sobre todo el dominio»?

\textbf{Respuesta E03.2:} \emph{Escribir aquí.}

\Needspace{4\baselineskip}

\textbf{E03.3, 3 puntos.} En E.12, indicar cuál factor vale cero y
explicar por qué eso no implica \(d\bar F/dy=0\).

\textbf{Respuesta E03.3:} \emph{Escribir aquí.}

\textbf{Condición esencial:} identificar perfil, dirección, tamaño y
pendiente, y justificar el término de borde. \textbf{Conexión:} A aplica
la envolvente; B distingue qué ocupaciones se fijan al derivar; C
obtiene fuerzas y bordes de una energía espacial; D reúne esas
ecuaciones. Fuente para nombre y definición:
\href{https://ocw.mit.edu/courses/18-325-topics-in-applied-mathematics-waves-and-imaging-fall-2015/834e98a77f64597a422c205b123cd6d7_MIT18_325F15_Appendix_A.pdf}{MIT
18.325, cálculo de variaciones, apéndice A}. Las expansiones y el
ejemplo de envolvente se calculan aquí.

\Needspace{12\baselineskip}

\section{E.6. De dos masas a una DOS ---
E05}\label{e.6.-de-dos-masas-a-una-dos-e05}

\textbf{Estado: ABIERTO. Meta:} separar forma de vibración, frecuencia,
conteo de modos y ocupación. DFPT se estudia después, en E08.

\subsection{E.6.1. Introducción desde modos
normales}\label{e.6.1.-introducciuxf3n-desde-modos-normales}

Dos masas iguales \(m\) pueden desplazarse a lo largo de una línea. Sus
desplazamientos son \(u_1,u_2\); tres resortes de constante \(K\) las
conectan entre sí y a paredes fijas. Hay dos coordenadas independientes,
luego dos modos normales. Un modo normal es un patrón de desplazamientos
que oscila con una frecuencia propia.

\begin{figure}
\centering
\includegraphics[width=0.98\linewidth,height=\textheight,keepaspectratio,alt={Figura E.6. Cada modo aporta un área de uno a la DOS. Las alturas de las líneas son una convención gráfica para mostrar esas áreas; no cuentan cuántos cuantos de vibración se han excitado.}]{../figuras/E05_modos_conteo.png}
\caption{Figura E.6. Cada modo aporta un área de uno a la DOS. Las
alturas de las líneas son una convención gráfica para mostrar esas
áreas; no cuentan cuántos cuantos de vibración se han excitado.}
\end{figure}

\subsection{E.6.2. Ejemplo resuelto: de fuerzas a un
conteo}\label{e.6.2.-ejemplo-resuelto-de-fuerzas-a-un-conteo}

La energía potencial suma el alargamiento al cuadrado de cada resorte:

\[
V=\frac K2[u_1^2+(u_2-u_1)^2+u_2^2],\qquad
m\ddot{\mathbf u}=-K\begin{pmatrix}2&-1\\-1&2\end{pmatrix}\mathbf u.
\tag{E.13}
\]

\(\mathbf u=(u_1,u_2)^T\) es la columna de desplazamientos; los puntos
indican derivadas temporales. La matriz aparece al derivar
\(-\partial V/\partial u_i\). Ensayar una oscilación
\(\mathbf u(t)=\mathbf v\cos(\omega t)\) convierte la ecuación de
movimiento en un problema de valores propios. Los vectores de forma
\(\mathbf v_1=(1,1)^T/\sqrt2\) y \(\mathbf v_2=(1,-1)^T/\sqrt2\) tienen
valores propios adimensionales 1 y 3:

\[
\omega_1=\sqrt{K/m},\qquad \omega_2=\sqrt{3K/m}.
\tag{E.14}
\]

\(\omega\) es frecuencia angular, medida en radianes por segundo. La
DOS, densidad de estados de vibración, reúne esas frecuencias en una
función de conteo. Para este sistema finito se escribe

\[
F_\omega(\omega)=\delta(\omega-\omega_1)+\delta(\omega-\omega_2),
\qquad \int_0^\infty F_\omega(\omega)d\omega=2.
\tag{E.15}
\]

La delta de Dirac \(\delta\) representa una línea de área unitaria
concentrada en una frecuencia. Ése es el motivo de sumar deltas: cada
modo debe contar una vez. No hace falta interpretar su altura puntual.
Una gráfica con picos anchos debe conservar sus áreas.

La ocupación \(n_j\) indica cuántos cuantos hay en el modo \(j\). Cada
cuanto tiene energía \(\hbar\omega_j\), donde \(\hbar=h/(2\pi)\) y \(h\)
es la constante de Planck. Con \(n_1=2,n_2=0\), la energía de
excitaciones es \(2\hbar\sqrt{K/m}\), omitiendo la energía de punto cero
de la referencia. La integral de la DOS sigue siendo 2. \textbf{Excitar
un modo no crea otro modo disponible.}

\subsection{E.6.3. Cambiar el eje sin cambiar el
conteo}\label{e.6.3.-cambiar-el-eje-sin-cambiar-el-conteo}

Si \(\Omega=\hbar\omega\) es la energía, la cantidad de modos en el
mismo intervalo debe cumplir \(F_\Omega d\Omega=F_\omega d\omega\). Como
\(d\Omega=\hbar d\omega\),

\[
F_\Omega(\Omega)=\frac1\hbar F_\omega(\Omega/\hbar).
\tag{E.16}
\]

No se copia la altura al cambiar el eje. Para frecuencia ordinaria
\(\nu\) en ciclos por segundo, se usa \(\Omega=h\nu\). Por eso un
archivo en THz se convierte con \(h\times10^{12}\), no con
\(\hbar\times10^{12}\). Una DOS por átomo, por celda o por volumen
requiere además indicar por qué cantidad se dividió el conteo. En tres
dimensiones hay tres coordenadas por átomo, pero ese conteo no se impone
a una tabla de procedencia desconocida.

\subsection{E.6.4. Actividad}\label{e.6.4.-actividad}

\Needspace{4\baselineskip}

\textbf{E05.1, 4 puntos.} Tres masas con desplazamiento escalar y
extremos fijos tienen tres modos. Escribir una DOS simbólica usando
frecuencias \(\omega_1,\omega_2,\omega_3\) e indicar su integral. No
calcular esas frecuencias.

\textbf{Respuesta E05.1:} \emph{Escribir aquí.}

\Needspace{4\baselineskip}

\textbf{E05.2, 3 puntos.} Si se duplica la ocupación de un modo, ¿qué
cambia en la energía y qué cambia en la integral de la DOS?

\textbf{Respuesta E05.2:} \emph{Escribir aquí.}

\Needspace{4\baselineskip}

\textbf{E05.3, 3 puntos.} Explicar por qué cambiar de frecuencia a
energía en el eje horizontal exige cambiar también la altura de la DOS.

\textbf{Respuesta E05.3:} \emph{Escribir aquí.}

\textbf{Condición esencial:} conservar el conteo y no confundirlo con
ocupación. \textbf{Conexión:} A distingue espectros electrónicos y
materiales; B necesita una DOS fonónica con unidades y normalización
verificadas; C recibe sus tasas; D exige esa comprobación antes de usar
datos. El ejemplo mecánico se deriva aquí. Para el contexto de modos de
un cristal: \href{https://arxiv.org/abs/cond-mat/0012092}{Baroni et al.,
dinámica de red y DFPT}, sección II.

\Needspace{12\baselineskip}

\section{E.7. Qué aporta DFPT al detector ---
E08}\label{e.7.-quuxe9-aporta-dfpt-al-detector-e08}

\textbf{Estado: ABIERTO. Meta:} explicar de dónde salen frecuencias y
pesos de interacción, sin resolver un programa de estructura
electrónica.

\subsection{E.7.1. Introducción: quién determina los
resortes}\label{e.7.1.-introducciuxf3n-quiuxe9n-determina-los-resortes}

En E05 se entregó \(K\) como dato. En un sólido, desplazar un átomo
modifica la distribución electrónica y las fuerzas sobre los demás.
\textbf{DFT}, teoría del funcional de la densidad, calcula el estado
electrónico de equilibrio mediante la densidad electrónica dentro de una
aproximación elegida. \textbf{DFPT}, teoría de perturbaciones del
funcional de la densidad, calcula cómo responde ese equilibrio a una
perturbación pequeña, por ejemplo un desplazamiento atómico.

«Autoconsistente» significa que la densidad determina el potencial
electrónico y que los estados calculados con ese potencial deben
devolver la misma densidad. No es simplemente mover un resorte
manteniendo congelada toda la nube electrónica.

\begin{figure}
\centering
\includegraphics[width=0.98\linewidth,height=\textheight,keepaspectratio,alt={Figura E.7. La curvatura de la energía determina fuerzas restauradoras. Los pesos de interacción responden a otra pregunta: cuánto acopla cada modo con los electrones. Ambos datos se necesitan para construir tasas.}]{../figuras/E08_curvatura_acoplamiento.png}
\caption{Figura E.7. La curvatura de la energía determina fuerzas
restauradoras. Los pesos de interacción responden a otra pregunta:
cuánto acopla cada modo con los electrones. Ambos datos se necesitan
para construir tasas.}
\end{figure}

\subsection{E.7.2. Ejemplo resuelto: curvatura y peso no son lo
mismo}\label{e.7.2.-ejemplo-resuelto-curvatura-y-peso-no-son-lo-mismo}

Primero, para un desplazamiento pequeño \(u\), expandir la energía de
equilibrio alrededor de \(u=0\) da \(E(u)=E_0+Ku^2/2+\cdots\). No hay
término lineal porque la fuerza inicial de equilibrio es cero. El
coeficiente \(K=E''(0)\) es la curvatura, con unidades de energía por
longitud al cuadrado.

Para varias coordenadas, \(u_I\) etiqueta un átomo y una dirección.
\(\Phi_{IJ}\) es una rigidez: la fuerza restauradora cumple
\(F_I\simeq-\sum_J\Phi_{IJ}u_J\). El signo menos proviene de derivar la
energía con \(F_I=-\partial E/\partial u_I\). La expansión de Taylor y
la matriz dinámica son

\[
E(\mathbf u)=E_0+\frac12\sum_{I,J}\Phi_{IJ}u_Iu_J+\cdots,
\quad \Phi_{IJ}=\frac{\partial^2E}{\partial u_I\partial u_J},
\quad \mathsf D_{IJ}=\frac{\Phi_{IJ}}{\sqrt{M_IM_J}}.
\tag{E.17}
\]

\(M_I\) es la masa del átomo de la coordenada \(I\). Dividir por las
masas convierte rigidez en frecuencia al cuadrado, como \(K/m\) en E05.
Los valores propios de \(\mathsf D\) son \(\omega^2\). En un cristal se
organiza además la dependencia espacial por vector de onda; aquí no se
requiere su álgebra.

Segundo, los elementos de interacción electrón--fonón, usualmente
llamados \(g\), miden cuánto conecta la perturbación del potencial a dos
modos electrónicos. Sus probabilidades de transición llevan pesos
proporcionales a \(|g|^2\). La función espectral \(\alpha^2F(\Omega)\)
reúne frecuencias \textbf{ponderadas por esa interacción}, no sólo el
conteo \(F(\Omega)\) de E05.

Como ejemplo de dos modos, definir sus energías \(\Omega_1=1\),
\(\Omega_2=2\) y pesos \(W_1=0.1\), \(W_2=0.3\), todos en una misma
unidad de energía. En la convención de eje energético con \(\alpha^2F\)
adimensional,

\[
\alpha^2F(\Omega)=W_1\delta(\Omega-\Omega_1)+W_2\delta(\Omega-\Omega_2),
\qquad
\lambda=2\int_0^\infty\frac{\alpha^2F(\Omega)}{\Omega}d\Omega.
\tag{E.18}
\]

\(\lambda\) es un parámetro adimensional de acoplamiento. La delta tiene
unidades inversas de energía; por eso \(W_j\) debe tener unidades de
energía para que \(\alpha^2F\) sea adimensional. Integrar cada delta
evalúa \(1/\Omega\) en su modo:

\[
\lambda=2\left(\frac{0.1}{1}+\frac{0.3}{2}\right)=0.5.
\tag{E.19}
\]

Duplicar ambos pesos duplica \(\lambda\) y no cambia el número de modos.
Estos valores son pedagógicos. No se ajusta con ellos NbN ni se deduce
automáticamente una temperatura crítica.

\subsection{E.7.3. La entrada concreta de
B}\label{e.7.3.-la-entrada-concreta-de-b}

Para evitar una normalización oculta, llamar \(g_{\rm ph}(\Omega)\) a la
densidad de modos por volumen y por energía. Si \(n(\Omega)\) es la
ocupación media de cada modo, la energía fonónica por volumen es

\[
u_{\rm ph}=\int_0^\infty g_{\rm ph}(\Omega)\,\Omega\,n(\Omega)d\Omega.
\tag{E.20}
\]

Se lee «modos por volumen e intervalo» por «energía de cada cuanto» por
«cuantos por modo», integrado sobre todos los intervalos. DFPT ayuda a
obtener modos y pesos materiales; B añade ocupaciones y factores
superconductores para formar tasas. DFPT por sí sola no entrega la
burbuja posterior al fotón ni el pulso del circuito. La revisión 0.4
comprueba la admisibilidad de los datos antes de atribuirles una
normalización absoluta.

\subsection{E.7.4. Actividad}\label{e.7.4.-actividad}

\Needspace{4\baselineskip}

\textbf{E08.1, 4 puntos.} Para dos modos con \(\Omega_1=1\),
\(\Omega_2=3\) y \(W_1=W_2=0.15\) en la misma unidad de energía,
calcular \(\lambda\) usando E.18.

\textbf{Respuesta E08.1:} \emph{Escribir aquí.}

\Needspace{4\baselineskip}

\textbf{E08.2, 3 puntos.} Dos espectros tienen el mismo número de modos
en cada frecuencia y la misma normalización, pero distintos pesos de
interacción. ¿Deben tener la misma DOS? ¿Deben tener el mismo
\(\alpha^2F\)? Explicar cada respuesta por separado.

\textbf{Respuesta E08.2:} \emph{Escribir aquí.}

\Needspace{4\baselineskip}

\textbf{E08.3, 3 puntos.} Un archivo dice «DOS» sin precisar si es por
átomo o por celda. ¿Basta multiplicarlo por una densidad de iones para
obtener \(u_{\rm ph}\)? Indicar el dato que falta.

\textbf{Respuesta E08.3:} \emph{Escribir aquí.}

\textbf{Condición esencial:} distinguir conteo, interacción y ocupación,
y comprobar sus unidades. \textbf{Conexión:} A define \(\lambda\) y el
alcance microscópico; B utiliza \(g_{\rm ph}\) y \(\alpha^2F\) en
balances y colisiones; C recibe sus intercambios y D registra las
entradas exigidas. Fuentes:
\href{https://arxiv.org/abs/cond-mat/0012092}{Baroni et al., DFPT} y
\href{https://www.quantum-espresso.org/Doc/ph_user_guide/node19.html}{Quantum
ESPRESSO, coeficientes electrón--fonón, ecuaciones (1)--(5)}. Se usa un
ejemplo espectral construido y no un cálculo DFPT nuevo.

\section{E.8. Respuestas, pauta y próxima
revisión}\label{e.8.-respuestas-pauta-y-pruxf3xima-revisiuxf3n}

Hay 21 respuestas pendientes, tres por clase activa. No se conserva una
pauta de E02 porque esa clase fue descartada. E03, E05 y las dos clases
nuevas de ocupaciones/Nambu se evaluarán con los requisitos reducidos de
esta versión. No se exige resolver actividades retiradas.

\textbf{Comentarios de aprendizaje:} \emph{Indicar qué símbolo o salto
dejó de entenderse y en qué paso. Si una clase quedó sin responder, se
revisará su dificultad antes de insistir con el mismo ejercicio.}

\textbf{Pauta de actividades:} pendiente de respuestas.
\textbf{Calificaciones:} sin asignar. En la próxima revisión se
conservará la evidencia de aprendizaje y se enlazarán los nuevos
conceptos con las clases efectivamente cerradas.

\end{document}
